\ReproPage{19. Companion videos: statistical fields and actual motion}
The random-vibration figures in this PDF are statistical, one-sigma RMS fields. An RMS field does not specify instantaneous displacement signs, relative phase, or a unique vibration history. The earlier camera-rotation videos preserve those fields and move only the viewpoint. Their captions identify that interpretation.

The additional fixed-camera motion videos use new native MAPDL transient calculations for the same complete assembly and the same three separate base-excitation directions. They show one reproducible synthetic realization of the specified environment. Their contours are instantaneous signed displacement, not RMS. Components, batteries, nickel strips, PCB, supports and fasteners remain in the solved assembly; a PCB-scoped picture restricts the displayed contour, not the calculation.

\ReportHeading{Inherited physical model}
Retain the 38 real modes from 485.964268212 to 3986.07476348 Hz, the locked six-bolt preload of 750 N per bolt, the original supports and the authorized nickel-only battery connections. Do not add battery-to-PCB bonding. Apply 2\% critical modal damping to the forced response. The eigensolution remains undamped. This is a linearized calculation about the retained preload/contact state: it does not update contact opening, sliding, impact, bolt loosening or plasticity during the random motion.

\ReportHeading{Time and amplitude shown in each clip}
Each displacement clip contains 480 actual result states, from 1.00006103515625 to 1.029296875 s, spaced by $1/16384$ s. Encoding those states at 60 frames/s gives an eight-second video and a playback slowdown of $16384/60=273.0667$. The displayed displacement vector is magnified by a constant factor of 50. The numerical contour values retain physical metres. The footer displays the actual simulated time, magnification and slowdown. The camera and color bounds remain fixed throughout each clip. The sequence runs forward once; it is not reversed to manufacture a repeating history.

The output is 3840 by 2160 pixels, H.264 in MP4, with no audio. These files are separate companions, not embedded PDF media. The sibling \path{videos/power_board_random_motion_4k60/index.html} is the motion-video index; \path{videos/power_board_4k60/index.html} contains the earlier camera views and modal companions. File availability is recorded in each index and its acceptance manifest.

\ReproPage{20. Reproduce the representative transient calculations}
Use the original calculation package \path{outputs/random_motion_20261005_2143/}, denoted \texttt{T} below. The report's \path{reproduction/motion/} contains the small scripts, input decks and audit records; retain the original binary solver dependencies separately. Adapt absolute host/Windows paths before replaying a script. Use fresh runtime directories and preserve the accepted source files.
\begin{enumerate}
\item Run \texttt{generate\_input.py} with NumPy. Use the four PSD points in this report, log--log interpolation, zero input outside 20--2000 Hz, and the original conversion $g_0=9.806$ m/s$^2$. Generate an eight-second periodic record at 65536 samples/s. At each positive Fourier frequency, use coefficient magnitude $N\sqrt{S_a(f)\,\Delta f/2}$ and an independent uniform random phase. Use seeds 20261005, 20261006 and 20261007 for X, Y and Z. Preserve the exact arrays and acceleration tables. Apply a half-cosine startup ramp over the first 0.1 s and extend the record to nine seconds.
\item Stage the matching modal database and mode, full-matrix, element and saved-state files, including all twelve distributed ranks, into a new Windows-local directory. Run the supplied no-SOLVE preflight, check the node/element inventory, 781 blocked mounting nodes and 38 retained frequencies, and seal the staged files with hashes. Do not run inside Workbench's generated source directory.
\item Restart the modal basis with \texttt{ANTYPE,MODAL,RESTART} and \texttt{MODCONT,,ON}. At the already blocked supports, identify enforced bases UX=1, UY=2 and UZ=3. Clear prior applied nodal forces and global acceleration. Generate the three enforced-motion vectors while preserving the eigensolution. Retain the resulting \texttt{.enf} files and all matching mode/restart files. Basis preparation used twelve distributed ranks; production transient integration used SMP with one rank.
\item For each separate excitation direction, stage a fresh transient runtime and perform its native preflight. Define \texttt{ANTYPE,TRANS}, \texttt{TRNOPT,MSUP,38,,1,YES,NMK,YES}, \texttt{DMPRAT,0.02}, constant $\Delta t=1/65536$ s, average-acceleration Newmark integration without added algorithmic damping, and \texttt{MCFOPT,1,0,0}. First solve the zero-input initial state. Read the acceleration table and apply \texttt{DVAL,base,ACC,\%PBACC\%,,OFF}; the last argument requests motion relative to the mounting base. Keep the other bases at zero. Solve to nine seconds.
\item Inspect the complete native log before accepting a run. Preserve 589825 time samples, \texttt{.mcf} and \texttt{.rdsp} histories, exact input, seed, runtime inventory, hashes and acceptance JSON. Reject missing, nonfinite, nonmonotone or incomplete histories. Failed trial runs remain diagnostic evidence and must not replace accepted data.
\item Use \texttt{verify\_transient.py} to compare the native modal coordinates with an independent Newmark integration. Discard the first second for steady RMS comparisons. Compare the nine original directional RMS hotspot values and require zero relative support displacement. The validation figures are listed on the next page.
\end{enumerate}

\ReproPage{21. Validate and expand the actual displacement states}
\ReportHeading{Accepted transient checks}
All three nine-second production histories completed with zero native solver errors. Relative support displacement remained zero. The coordinate checks and largest directional RMS differences for each excitation were:
\begin{center}\begin{tabular}{c r r}\toprule
Base & Relative $L_2$ coordinate error & Largest absolute RMS difference\\\midrule
X & $1.2568\times10^{-5}$ & 0.0545\%\\
Y & $3.7921\times10^{-6}$ & 0.1577\%\\
Z & $1.2407\times10^{-5}$ & 0.0334\%\\\bottomrule\end{tabular}\end{center}
These are numerical cross-checks against independent integration and the earlier PSD response; they are not physical validation of the inherited material, contact or damping assumptions. The retained modal effective-mass fractions remain those reported earlier; the video calculation does not establish modal or mesh convergence.

\begin{enumerate}
\item Copy each accepted transient state into a fresh expansion directory. Resume its matching database and request a nodal expansion of the selected stationary window. The accepted native commands include \texttt{EXPASS,ON}, \texttt{NUMEXP,480,1,1.029296875,NO}, \texttt{OUTRES,ALL,NONE}, \texttt{OUTRES,NSOL,ALL}, then \texttt{SOLVE}. These expansions contain nodal displacements; they do not contain element stresses.
\item Archive and hash each completed \texttt{file.rst} outside the Workbench-generated analysis directory. The final file contains 480 states beginning at 1.00006103515625 s. Verify the actual time list; do not infer the saved times from the requested start alone.
\item In Mechanical, create a new Transient Structural branch and link its modal initial condition to the original 38-mode analysis (original object 4254). Set the end time to 9 s and the time step to $1/65536$ s \emph{before} importing the result by reference. A default time step larger than the imported window can prevent evaluation despite a valid result file.
\item Read the immutable archived result with \texttt{ReadGivenAnsysResultFileByReference}, using MKS units. Create the required directional displacement result with \texttt{By=ResultSet}. For PCB-only contours, scope it to the original PCB body 42342 after verifying that body's identity. Reimports and regenerated models can change object IDs; do not copy IDs without checking their names, ownership and geometry.
\item Export the native displacement vector at seven validation nodes for sets 1, 2, 120, 240, 360 and 480. Compare with $\mathbf{u}(t)=\sum_{j=1}^{38}\boldsymbol{\phi}_j q_j(t)$, using the original modal shapes and native modal coordinates. The accepted relative $L_2$ errors are $6.59\times10^{-8}$ for X, $1.75\times10^{-7}$ for Y and $8.85\times10^{-8}$ for Z. Require the 480 times to increase by exactly $1/16384$ s within the recorded floating-point tolerance. Preserve these checks in \texttt{expand\_X/Y/Z/ACCEPTED.json}.
\end{enumerate}

\ReproPage{22. Render, encode and check the motion videos}
\begin{enumerate}
\item Evaluate the chosen directional result over all 480 states to obtain the largest absolute value in the window. Choose fixed symmetric legend bounds that include that value, with ten color bands. Keep those bounds unchanged at every frame. Show all unsuppressed assembly bodies for whole-assembly views; use the verified PCB result for a scoped view. Hide mesh edges, date/time annotations and min/max arrows.
\item Set true deformation scaling with multiplier 50. For the oblique battery-side view, use ViewVector $(1,-1,-1)$ and UpVector $(-0.3,1,-1.3)$. Fit once, then increase SceneHeight by 1.3. For the scoped PCB views, use the specified isometric orientation. Do not refit or move the camera during playback.
\item The accepted explicit-frame method selects each actual result set, evaluates it, activates it, reasserts the fixed color bounds and calls native Mechanical \texttt{ExportImage}. Use 3840 by 2160 pixels, \texttt{CurrentGraphicsDisplay=False}, and font magnification 0.85. Use \texttt{CalculateTimeHistory=False} for this explicit per-set method: every frame is independently evaluated at its own set. Record set number, actual time, maximum and minimum next to each PNG. A four-frame proof used sets 1, 2, 120 and 480 before production.
\item The earlier accepted X/Z exports used native \texttt{ExportAnimation}. That route requires \texttt{CalculateTimeHistory=True}, RangeType=ResultSets, ForwardBackwardMode=0, 480 frames and an explicit native \texttt{SetExportStreamWidthHeight(3840,2160)}. The public width/height settings alone previously produced viewport-sized frames. \texttt{CalculateTimeHistory=False} with this animation route animates a single selected state and is not an acceptable history. Its end-of-encoding hang motivated the explicit-frame method above.
\item Render to a Windows-local directory, then verify and archive the native PNGs before transfer. Check exactly 480 complete, distinct images at 3840 by 2160, correct forward time order, fixed camera and bounds, and unmodified model pixels. Preserve SHA-256 hashes and the per-frame audit. A completed MP4 or a GUI progress bar alone is not sufficient evidence.
\item Use the supplied hardware-encoding script on the host. It uses \texttt{h264\_videotoolbox}, software fallback disabled, 60 fps, High profile/Level 5.2, YUV420p, an 80 Mbit/s target and 120 Mbit/s maximum, BT.709 metadata, and fast-start MP4. Add clear English titles, physical time, magnification and slowdown. A readable fixed legend uses the verified native ten-color palette and exact fixed bounds; it does not change the native geometry or contour colors.
\item Check the complete MP4 with \texttt{ffprobe}: 3840 by 2160, rate 60/1, 480 frames and eight seconds. Decode the entire file with \texttt{ffmpeg -v error -i video.mp4 -f null -}. Inspect representative frames and playback for clipped geometry, stale min/max labels, obscured legends, camera motion, incorrect result scope or repeated/interpolated states. Publish only accepted files in the video index.
\end{enumerate}
The Windows VM uses the Parallels WDDM graphics device; encoding uses the Mac's hardware video engine. The physical GPU is not directly passed through to Windows. Keep each clip's manifest and encoding command with the report reproduction package. Stress histories require a separately validated tensor-recovery procedure; displacement-only expansion files cannot support a stress-motion video.
